\worksheettitle
  {M06-R. Идём по границе}
  {\modulemark{M06-R} Коррекционная карточка}

\studentnameblock

\begin{warningbox}[Назовите ошибку]
  Буквы в названии нельзя переставлять произвольно. Сначала проверьте путь по
  сторонам, затем соседство концов каждого отрезка.
\end{warningbox}

\begin{practicebox}[Шаг 1. Восстановите обход]
  Для пятиугольника $ABCDE$ продолжите имена:
  \[
    A-B-\answerblank[12mm]-\answerblank[12mm]-\answerblank[12mm],
  \]
  \[
    A-\answerblank[12mm]-\answerblank[12mm]-\answerblank[12mm]
      -\answerblank[12mm].
  \]
\end{practicebox}

\begin{practicebox}[Шаг 2. Проверьте соседство]
  В пятиугольнике $ABCDE$:
  \begin{enumerate}[label=\alph*)]
    \item с $A$ соседние вершины \answerblank[15mm] и \answerblank[15mm];
    \item $AB$ --- \answerblank[22mm];
    \item $AC$ --- \answerblank[22mm].
  \end{enumerate}
\end{practicebox}

\begin{practicebox}[Шаг 3. Исправьте название]
  Запись $ACDEB$ неверна.

  Первый нарушенный переход: \answerblank[25mm]

  Два верных имени: \answerblank[42mm], \answerblank[42mm].
\end{practicebox}

\begin{practicebox}[Шаг 4. Исправьте объяснение]
  Ученик сказал: «$AE$ --- диагональ, потому что буквы не стоят рядом
  в алфавите».

  Причина: \answerblank[95mm]

  Исправление: \answerblank[95mm]
\end{practicebox}

\begin{checkpointbox}[Повторная проверка]
  Для шестиугольника $KLMNPQ$:
  \begin{enumerate}[label=\textbf{\arabic*.}]
    \item запишите два допустимых имени;
    \item определите, чем являются $KL$, $KN$ и $QK$;
    \item назовите все диагонали из $K$.
  \end{enumerate}
  \writinglines{3}
\end{checkpointbox}
